% Nouveau EGA V (cote 25): archive cote 25, site page 6 (PDF page 7).
% typescript. Draft compared: none.
% AI-transcribed by Claude Opus 5.5 (claude-opus-5-5, adaptive thinking, effort high), Message Batches API, batch ega25b (msgbatch_01A7n5YuXubUPNK1zsepi8RG), 2026-11-13; not yet checked by eye.
\scanpage{25/6}
\mspage{3}

Faisons maintenant le changement de base $S'=\overset{\hand{\vee}}{P}\longrightarrow S$, et considérons
la section diagonale (ou\hand{``}section générique\hand{''}) de $\overset{\hand{\vee}}{P}_{S'} = P\,(\overset{\hand{\vee}}{\underline{E}}_{S'})$ sur $S'$ :
on trouve un sous-schéma fermé $\underline{H}$ de $P_{S'} = P\times_S\overset{\hand{\vee}}{P}$ , appelé parfois
\uline{schéma d'incidence} entre $P$ et $\overset{\hand{\vee}}{P}$ , défini par l'Idéal image de
l'homomorphisme canonique
\[ \underline{O}_P(-1)\otimes_S\underline{O}_{\overset{\hand{\vee}}{P}}(-1)\longrightarrow \underline{O}_{P\times_S\overset{\hand{\vee}}{P}} \qquad ; \]
d'après ce qu'on a vu, c'est un fibré projectif sur $\overset{\hand{\vee}}{P}$, et par raison de
symétrie, c'est également un fibré projectif sur $P$. Bien entendu, on
retrouve les\deleted{\xout{\unsure{définitions des}}}\inserted{hyperplans ``spéciaux''} $H_{\hand{\xi}}$ ( pour $\hand{\xi}$ section de $\overset{\hand{\vee}}{P}$ sur $S$ )\hand{,} à partir
du "hyperplan universel" $\underline{H}$\hand{,} en prenant son image inverse par l\over{a}{e} change-%
ment de base $S\xrightarrow{\hand{\xi}}\overset{\hand{\vee}}{P}$ . La même remarque s'applique pour tout point
$\hand{\xi}$ de $\overset{\hand{\vee}}{P}$ à valeurs dans un $S$-préschéma quelconque $S'$, qui (considéré comme
section de $\overset{\hand{\vee}}{P}_{S'}$ sur $S'$) permet de définir un $H_{\hand{\xi}}\subset P_{S'}$ : ce dernier n'est
autre que l'image inverse de $\underline{H}$ par le changement de base $S'\xrightarrow{\hand{\xi}}\overset{\hand{\vee}}{P}$.

Nous supposons par la suite donné un préschéma de type fini $X$ sur $S$,
et un $S$-morphisme
\[ f\colon X\longrightarrow P \quad . \]
\margin{$|$}Un des principaux objets\hand{,} du présent paragraphe est d'étudier, pour tout
\margin{$|$}hyperplan $H_{\hand{\xi}}$ de $P$, \inserted{$(\xi\in\overset{\vee}{P}(S))$} son image inverse
\[ Y_{\hand{\xi}} = f^{-1}(H_{\hand{\xi}}) = X\times_P H_{\hand{\xi}} \quad , \]
et notamment de mettre en relations le propriétés de $X$ et des $Y_{\hand{\xi}}$ . Com-%
me d'habitude, on aura également à considérer des $\overset{\hand{\vee}}{P}(S')$, $S'$ \xout{\illegible{1}}
un $S$-préschéma quelconque, dans ce cas $H_{\hand{\xi}}$ est un hyperplan de $P_{S'}$ \hand{,}
et on posera encore
\[ Y_{\hand{\xi}} = f_{S'}{}^{-1}(H_{\hand{\xi}}) = X_{S'}\times_{P_{S'}}H_{\hand{\xi}} = X\times_P H_{\hand{\xi}} \quad , \]
où l'indice $S'$ désigne comme d'\unsure{habitude|hhabxitude} \xout{\illegible{1}} l'effet du changement de base
$S'\longrightarrow S$, et où dans le dernier terme on considère $H_{\hand{\xi}}$ comme $P$-préschéma.

\mspage{6}
